% Fichier engendré par outils/generer-tex.py à partir de 08-geometrie-vecteurs-produit-scalaire-espace.
% Les démonstrations en français sont transcrites du fichier Lean (skill transcrire-preuve-lean).
\documentclass[11pt,a4paper]{article}

% Compile aussi bien avec pdflatex qu'avec xelatex ou tectonic.
\usepackage{iftex}
\ifPDFTeX
  \usepackage[T1]{fontenc}
  \usepackage[utf8]{inputenc}
\else
  \usepackage{fontspec}
\fi
\usepackage[french]{babel}
\usepackage{amsmath,amssymb,amsthm}
\usepackage[margin=2.5cm]{geometry}
\usepackage[hidelinks]{hyperref}

\theoremstyle{definition}
\newtheorem{definition}{Définition}
\theoremstyle{plain}
\newtheorem{theoreme}{Théorème}
\newtheorem{lemme}[theoreme]{Lemme}

% Renvoi à la déclaration Lean, une ligne après la démonstration, aligné à droite.
\newcommand{\source}[2]{\par\smallskip\noindent\hfill{\footnotesize\href{#1}{\texttt{#2}}}\par\smallskip}

\title{Géométrie : vecteurs, produit scalaire, espace}
\date{}

\begin{document}
\maketitle
% Les identifiants Lean en machine à écrire ne se coupent pas : on tolère des
% espacements plus lâches plutôt que des lignes qui débordent.
\sloppy

\section{GeometrieVecteursProduitScalaireEspace}

\noindent
Lycée \textemdash{} section \og{} Géométrie : vecteurs, produit scalaire, espace \fg{}, de
la seconde à la terminale S.

\medskip\noindent
Deux cadres se partagent le chapitre. Les énoncés qui parlent de coordonnées vivent dans
le plan repéré $\mathbb{R}^{2}$ ou l'espace repéré $\mathbb{R}^{3}$ munis de la distance
euclidienne : un point y \emph{est} la donnée de ses coordonnées, notées $a_0$, $a_1$,
$a_2$. Les autres \textemdash{} produit scalaire, Al-Kashi, médiane \textemdash{} sont
écrits dans un espace préhilbertien quelconque, c'est-à-dire sans repère, comme le fait le
programme lorsqu'il raisonne sur des vecteurs.

\medskip\noindent
Le produit scalaire est noté $\langle u, v\rangle$ ; c'est celui de la bibliothèque, dont
la version réelle est symétrique et bilinéaire.

\medskip\noindent
\textbf{Trois énoncés du programme ne sont pas entièrement traités} et sont signalés comme
tels à leur place : la classification des positions relatives de droites et de plans, le
théorème du toit, et les intersections droite-plan et plan-plan.

\subsection{Vecteurs : Chasles et colinéarité}

\begin{theoreme}
Relation de Chasles :
\[ \overrightarrow{AB} + \overrightarrow{BC} = \overrightarrow{AC} . \]
\end{theoreme}

\begin{proof}
Écrite avec les vecteurs $\overrightarrow{AB} = b - a$, la relation devient
\[ (b - a) + (c - b) = c - a , \]
où $b$ et $-b$ s'éliminent : c'est un simple calcul dans un groupe abélien, et il n'y a
rien de plus dans la relation de Chasles.
\end{proof}
\source{https://github.com/Commutator-IO/learn-lean/blob/main/cours/02-lycee/08-geometrie-vecteurs-produit-scalaire-espace/GeometrieVecteursProduitScalaireEspace.lean\#L20}{GeometrieVecteursProduitScalaireEspace.lean\#L20}

\begin{theoreme}
Deux vecteurs sont colinéaires si et seulement si l'un appartient à la droite vectorielle
engendrée par l'autre :
\[ (\exists k \in \mathbb{R},\ v = k u) \iff v \in \mathrm{Vect}(u) . \]
\end{theoreme}

\begin{proof}
La bibliothèque décrit l'espace engendré par un seul vecteur : $v$ y appartient si et
seulement s'il existe un scalaire $k$ tel que $ku = v$. Les deux membres de l'équivalence
disent donc la même chose, à l'ordre près de l'égalité, et chaque sens se démontre en
reprenant le $k$ fourni par l'autre.

\medskip\noindent
C'est la traduction vectorielle du parallélisme : deux droites $(AB)$ et $(CD)$ sont
parallèles exactement quand $\overrightarrow{AB}$ et $\overrightarrow{CD}$ engendrent la
même droite vectorielle.
\end{proof}
\source{https://github.com/Commutator-IO/learn-lean/blob/main/cours/02-lycee/08-geometrie-vecteurs-produit-scalaire-espace/GeometrieVecteursProduitScalaireEspace.lean\#L25}{GeometrieVecteursProduitScalaireEspace.lean\#L25}

\begin{theoreme}
Critère de colinéarité par le déterminant : pour $u \ne \vec 0$ de coordonnées
$(x, y)$ et $v$ de coordonnées $(x', y')$,
\[ u \text{ et } v \text{ colinéaires} \iff xy' - x'y = 0 . \]
\end{theoreme}

\begin{proof}
Comme $u$ n'est pas le vecteur nul, l'une au moins de ses coordonnées ne l'est pas : sinon
les deux coordonnées seraient nulles et $u$ le serait aussi.

\smallskip\noindent
\emph{Sens direct.} Si $v = ku$, alors $x' = kx$ et $y' = ky$, donc
\[ xy' - x'y = x(ky) - (kx)y = 0 . \]

\smallskip\noindent
\emph{Sens réciproque.} Supposons $xy' - x'y = 0$ et distinguons deux cas.

Si $x \ne 0$, il suffit de prendre $k = x'/x$. La première coordonnée est bien
$k x = x'$ ; pour la seconde, $ky = x'y/x$ vaut $y'$ précisément parce que
$xy' = x'y$.

Si $x = 0$, alors $y \ne 0$ et l'on prend $k = y'/y$ : la seconde coordonnée est correcte,
et la première l'est parce que le déterminant se réduit alors à $-x'y = 0$, d'où $x' = 0$
puisque $y \ne 0$.
\end{proof}
\source{https://github.com/Commutator-IO/learn-lean/blob/main/cours/02-lycee/08-geometrie-vecteurs-produit-scalaire-espace/GeometrieVecteursProduitScalaireEspace.lean\#L36}{GeometrieVecteursProduitScalaireEspace.lean\#L36}

\subsection{Système linéaire et base du plan}

\begin{theoreme}
Un système linéaire $2 \times 2$ de déterminant non nul a une solution et une seule :
si $ad - bc \ne 0$, le système
\[ \begin{cases} ax + by = e \\ cx + dy = f \end{cases} \]
a pour unique solution
\[ x = \frac{ed - bf}{ad - bc}, \qquad y = \frac{af - ec}{ad - bc} . \]
\end{theoreme}

\begin{proof}
\emph{Existence.} On remplace $x$ et $y$ par les valeurs annoncées et l'on réduit au même
dénominateur $ad - bc$, non nul par hypothèse :
\[ a\,\frac{ed - bf}{ad - bc} + b\,\frac{af - ec}{ad - bc}
   = \frac{aed - abf + abf - bec}{ad - bc}
   = \frac{e\,(ad - bc)}{ad - bc} = e , \]
et le même calcul donne la seconde équation.

\emph{Unicité.} Soit $(u, v)$ une solution. En combinant les deux équations \textemdash{}
la première multipliée par $d$, moins la seconde multipliée par $b$ \textemdash{} les
termes en $v$ disparaissent :
\[ u\,(ad - bc) = ed - bf , \]
et symétriquement, la seconde multipliée par $a$ moins la première par $c$ donne
\[ v\,(ad - bc) = af - ec . \]
Comme $ad - bc \ne 0$, on peut diviser, et $(u, v)$ est bien le couple annoncé.
\end{proof}
\source{https://github.com/Commutator-IO/learn-lean/blob/main/cours/02-lycee/08-geometrie-vecteurs-produit-scalaire-espace/GeometrieVecteursProduitScalaireEspace.lean\#L61}{GeometrieVecteursProduitScalaireEspace.lean\#L61}

\begin{theoreme}
Décomposition unique dans une base du plan : si $u$ et $v$ ont un déterminant non nul,
tout vecteur $w$ s'écrit d'une seule façon
\[ w = \alpha u + \beta v . \]
\end{theoreme}

\begin{proof}
Les deux coordonnées de l'égalité $w = \alpha u + \beta v$ forment exactement le système
\[ \begin{cases} \alpha\,u_0 + \beta\,v_0 = w_0 \\ \alpha\,u_1 + \beta\,v_1 = w_1 \end{cases} \]
dont le déterminant est $u_0 v_1 - v_0 u_1 \ne 0$. Le théorème précédent lui donne une
solution et une seule.

Pour l'existence, on vérifie l'égalité vectorielle coordonnée par coordonnée : c'est le
sens \og{} solution du système \fg{}. Pour l'unicité, on part d'une décomposition
quelconque, on lit ses deux coordonnées, et l'unicité du système conclut.
\end{proof}
\source{https://github.com/Commutator-IO/learn-lean/blob/main/cours/02-lycee/08-geometrie-vecteurs-produit-scalaire-espace/GeometrieVecteursProduitScalaireEspace.lean\#L84}{GeometrieVecteursProduitScalaireEspace.lean\#L84}

\subsection{Milieu et distance}

\begin{theoreme}
Les coordonnées du milieu sont les moyennes des coordonnées, et la distance se calcule par
le théorème de Pythagore :
\[ I_i = \frac{a_i + b_i}{2} , \qquad
   AB = \sqrt{(b_0 - a_0)^{2} + (b_1 - a_1)^{2}} . \]
\end{theoreme}

\begin{proof}
Le milieu est défini dans la bibliothèque par $\frac{1}{2}(a + b)$ ; sa $i$-ième
coordonnée est donc bien la demi-somme des coordonnées.

La distance euclidienne est définie comme la racine carrée de la somme des carrés des
écarts de coordonnées ; en développant la somme sur les deux indices et en remplaçant la
distance entre réels par la valeur absolue de la différence, dont le carré est le carré de
la différence, on obtient la formule annoncée.
\end{proof}
\source{https://github.com/Commutator-IO/learn-lean/blob/main/cours/02-lycee/08-geometrie-vecteurs-produit-scalaire-espace/GeometrieVecteursProduitScalaireEspace.lean\#L105}{GeometrieVecteursProduitScalaireEspace.lean\#L105}

\subsection{Équation de droite}

\begin{theoreme}
Une droite non verticale d'équation $y = mx + p$ a aussi pour équation cartésienne
$mx - y + p = 0$ ; et deux telles droites sont strictement parallèles si et seulement si
\[ m = m' \quad\text{et}\quad p \ne p' . \]
\end{theoreme}

\begin{proof}
La première équivalence est le passage d'un membre à l'autre dans une égalité de réels.

\smallskip\noindent
\emph{Parallélisme, sens direct.} Supposons les deux droites sans point commun.

Si l'on avait $m \ne m'$, le réel $x = \dfrac{p' - p}{m - m'}$ vérifierait
$mx + p = m'x + p'$ \textemdash{} il suffit de réduire au même dénominateur, licite
puisque $m - m' \ne 0$ \textemdash{} et les droites se couperaient. Donc $m = m'$.

Si l'on avait $p = p'$, les deux droites se couperaient en $x = 0$. Donc $p \ne p'$.

\smallskip\noindent
\emph{Sens réciproque.} Si $m = m'$ et $p \ne p'$, une abscisse commune donnerait
$mx + p = mx + p'$, donc $p = p'$ : contradiction. Aucun point n'est commun.
\end{proof}
\source{https://github.com/Commutator-IO/learn-lean/blob/main/cours/02-lycee/08-geometrie-vecteurs-produit-scalaire-espace/GeometrieVecteursProduitScalaireEspace.lean\#L122}{GeometrieVecteursProduitScalaireEspace.lean\#L122}

\subsection{Produit scalaire}

\begin{theoreme}
Trois expressions du produit scalaire :
\[ \langle u, v\rangle = \frac{\|u + v\|^{2} - \|u\|^{2} - \|v\|^{2}}{2} , \qquad
   \langle a, b\rangle = a_0 b_0 + a_1 b_1 , \]
\[ \langle u, v\rangle = \|u\|\,\|v\|\cos\bigl(\widehat{u, v}\bigr) . \]
\end{theoreme}

\begin{proof}
\emph{Par les normes.} L'identité $\|u + v\|^{2} = \|u\|^{2} + 2\langle u, v\rangle +
\|v\|^{2}$, démontrée dans la bibliothèque, se résout en $\langle u, v\rangle$ : c'est
l'identité de polarisation.

\emph{Par les coordonnées.} Dans le plan repéré, le produit scalaire est défini comme la
somme $\sum_i a_i b_i$, que l'on développe sur les deux indices.

\emph{Par le cosinus.} La bibliothèque \emph{définit} l'angle de deux vecteurs par
\[ \cos\bigl(\widehat{u, v}\bigr) = \frac{\langle u, v\rangle}{\|u\|\,\|v\|} , \]
et démontre l'identité $\cos(\widehat{u,v}) \|u\| \|v\| = \langle u, v\rangle$ y compris
dans les cas dégénérés où l'un des vecteurs est nul. La troisième expression n'est donc
pas une définition indépendante : c'est celle-là qui, dans Mathlib, définit l'angle, alors
que le lycée fait l'inverse.
\end{proof}
\source{https://github.com/Commutator-IO/learn-lean/blob/main/cours/02-lycee/08-geometrie-vecteurs-produit-scalaire-espace/GeometrieVecteursProduitScalaireEspace.lean\#L146}{GeometrieVecteursProduitScalaireEspace.lean\#L146}

\begin{theoreme}
Symétrie et bilinéarité :
\[ \langle u, v\rangle = \langle v, u\rangle, \qquad
   \langle u + v, w\rangle = \langle u, w\rangle + \langle v, w\rangle, \qquad
   \langle k u, w\rangle = k\,\langle u, w\rangle . \]
\end{theoreme}

\begin{proof}
Les trois propriétés sont celles de la bibliothèque, citées : sur un espace préhilbertien
\emph{réel}, le produit scalaire est symétrique et linéaire en chaque variable. Sur un
espace complexe la symétrie deviendrait une symétrie hermitienne, ce qui n'est pas le cadre
du programme.
\end{proof}
\source{https://github.com/Commutator-IO/learn-lean/blob/main/cours/02-lycee/08-geometrie-vecteurs-produit-scalaire-espace/GeometrieVecteursProduitScalaireEspace.lean\#L160}{GeometrieVecteursProduitScalaireEspace.lean\#L160}

\begin{theoreme}
Deux vecteurs sont orthogonaux si et seulement si les diagonales du parallélogramme qu'ils
construisent ont même longueur :
\[ \langle u, v\rangle = 0 \iff \|u + v\| = \|u - v\| . \]
\end{theoreme}

\begin{proof}
Les deux normes étant positives, l'égalité $\|u+v\| = \|u-v\|$ équivaut à l'égalité de
leurs carrés. Or
\[ \|u + v\|^{2} = \|u\|^{2} + 2\langle u, v\rangle + \|v\|^{2} , \qquad
   \|u - v\|^{2} = \|u\|^{2} - 2\langle u, v\rangle + \|v\|^{2} , \]
et la différence des deux vaut $4\langle u, v\rangle$. Les carrés sont donc égaux si et
seulement si le produit scalaire est nul.

\medskip\noindent
L'énoncé du programme, \og{} $u \perp v$ si et seulement si $u \cdot v = 0$ \fg{}, est en
Lean une définition et non un théorème : l'orthogonalité y \emph{est} la nullité du
produit scalaire. C'est pourquoi la transcription retient plutôt la caractérisation
géométrique, qui a un contenu \textemdash{} le rectangle est le parallélogramme dont les
diagonales sont égales.
\end{proof}
\source{https://github.com/Commutator-IO/learn-lean/blob/main/cours/02-lycee/08-geometrie-vecteurs-produit-scalaire-espace/GeometrieVecteursProduitScalaireEspace.lean\#L169}{GeometrieVecteursProduitScalaireEspace.lean\#L169}

\subsection{Al-Kashi, médiane, loi des sinus, aire}

\begin{theoreme}
Théorème d'Al-Kashi :
\[ AC^{2} = AB^{2} + BC^{2} - 2\,AB \cdot BC \cos\widehat{ABC} . \]
\end{theoreme}

\begin{proof}
Résultat de la bibliothèque, à ceci près qu'elle l'énonce avec des produits
$AC \times AC$ plutôt que des carrés, et avec la distance $CB$ là où le programme écrit
$BC$. La symétrie de la distance permet de passer de l'une à l'autre, et le reste est une
mise en forme.

L'argument du théorème est le développement de
$\|\overrightarrow{BA} - \overrightarrow{BC}\|^{2}$, où le double produit fait apparaître
$\langle \overrightarrow{BA}, \overrightarrow{BC}\rangle$, c'est-à-dire
$AB \cdot BC \cos\widehat{ABC}$. Le théorème de Pythagore en est le cas particulier
$\widehat{ABC} = \pi/2$.
\end{proof}
\source{https://github.com/Commutator-IO/learn-lean/blob/main/cours/02-lycee/08-geometrie-vecteurs-produit-scalaire-espace/GeometrieVecteursProduitScalaireEspace.lean\#L177}{GeometrieVecteursProduitScalaireEspace.lean\#L177}

\begin{theoreme}
Formule des trois normes (identité du parallélogramme) et théorème de la médiane :
\[ \|u + v\|^{2} + \|u - v\|^{2} = 2\bigl(\|u\|^{2} + \|v\|^{2}\bigr) , \]
\[ \|u\|^{2} + \|v\|^{2} = 2\,\Bigl\|\tfrac{1}{2}(u + v)\Bigr\|^{2}
   + \tfrac{1}{2}\,\|u - v\|^{2} . \]
\end{theoreme}

\begin{proof}
La première identité s'obtient en développant les deux carrés : les doubles produits
$+2\langle u, v\rangle$ et $-2\langle u, v\rangle$ s'éliminent, et il reste deux fois la
somme des carrés des normes.

La seconde est la même, réécrite. En sortant le facteur $\frac{1}{2}$ de la norme
\textemdash{} $\|\frac{1}{2}w\| = \frac{1}{2}\|w\|$ \textemdash{} le membre de droite
devient $\frac{1}{2}\|u+v\|^{2} + \frac{1}{2}\|u-v\|^{2}$, c'est-à-dire la moitié du membre
de gauche de la première identité, donc $\|u\|^{2} + \|v\|^{2}$.

\medskip\noindent
Sous cette forme, $\frac{1}{2}(u+v)$ est le vecteur qui joint le sommet au milieu du côté
opposé : c'est le théorème de la médiane du programme,
$MA^{2} + MB^{2} = 2\,MI^{2} + \frac{AB^{2}}{2}$.
\end{proof}
\source{https://github.com/Commutator-IO/learn-lean/blob/main/cours/02-lycee/08-geometrie-vecteurs-produit-scalaire-espace/GeometrieVecteursProduitScalaireEspace.lean\#L186}{GeometrieVecteursProduitScalaireEspace.lean\#L186}

\begin{theoreme}
Loi des sinus :
\[ \sin\widehat{ABC} \times BC = \sin\widehat{CAB} \times CA . \]
\end{theoreme}

\begin{proof}
Résultat de la bibliothèque, cité tel quel. Sous la forme du programme, en divisant par le
produit des deux longueurs, il s'écrit
\[ \frac{\sin \widehat{A}}{BC} = \frac{\sin \widehat{B}}{CA} , \]
le rapport commun étant l'inverse du diamètre du cercle circonscrit.
\end{proof}
\source{https://github.com/Commutator-IO/learn-lean/blob/main/cours/02-lycee/08-geometrie-vecteurs-produit-scalaire-espace/GeometrieVecteursProduitScalaireEspace.lean\#L199}{GeometrieVecteursProduitScalaireEspace.lean\#L199}

\begin{theoreme}
Aire d'un triangle : dans un plan orienté, la valeur absolue du déterminant de deux
vecteurs vaut
\[ |\det(u, v)| = \|u\|\,\|v\|\,\sin\bigl(\widehat{u, v}\bigr) , \]
de sorte que le triangle qu'ils portent a pour aire $\frac{1}{2}ab\sin C$.
\end{theoreme}

\begin{proof}
La bibliothèque démontre l'identité
\[ \langle u, v\rangle^{2} + \det(u, v)^{2} = \|u\|^{2}\,\|v\|^{2} , \]
qui est la forme algébrique de \og{} $\cos^{2} + \sin^{2} = 1$ \fg{} pour deux vecteurs.
En y remplaçant $\langle u, v\rangle$ par $\|u\|\|v\|\cos(\widehat{u,v})$, il vient
\[ \det(u, v)^{2}
   = \|u\|^{2}\|v\|^{2}\bigl(1 - \cos^{2}(\widehat{u,v})\bigr)
   = \bigl(\|u\|\,\|v\|\sin(\widehat{u,v})\bigr)^{2} . \]

Les deux membres de l'égalité cherchée sont positifs : à gauche une valeur absolue, à
droite un produit de normes par un sinus, positif parce que l'angle non orienté est compris
entre $0$ et $\pi$. Deux réels positifs de même carré étant égaux, on conclut.
\end{proof}
\source{https://github.com/Commutator-IO/learn-lean/blob/main/cours/02-lycee/08-geometrie-vecteurs-produit-scalaire-espace/GeometrieVecteursProduitScalaireEspace.lean\#L207}{GeometrieVecteursProduitScalaireEspace.lean\#L207}

\subsection{Cercle}

\begin{theoreme}
Équation cartésienne d'un cercle : pour $r \ge 0$,
\[ \Omega M = r \iff (x_M - x_\Omega)^{2} + (y_M - y_\Omega)^{2} = r^{2} . \]
\end{theoreme}

\begin{proof}
La distance est la racine carrée de la somme des carrés des écarts de coordonnées, somme
que l'on note $S$.

Si $\sqrt{S} = r$, alors $r^{2} = S$ en élevant au carré, ce qui est licite parce que $S$
est positive. Réciproquement, si $S = r^{2}$, alors $\sqrt{S} = \sqrt{r^{2}} = r$, la
dernière égalité utilisant $r \ge 0$.

\medskip\noindent
L'hypothèse $r \ge 0$ n'est pas décorative : sans elle l'équivalence tombe, un rayon
négatif ne pouvant être une distance alors que son carré, lui, reste positif.
\end{proof}
\source{https://github.com/Commutator-IO/learn-lean/blob/main/cours/02-lycee/08-geometrie-vecteurs-produit-scalaire-espace/GeometrieVecteursProduitScalaireEspace.lean\#L226}{GeometrieVecteursProduitScalaireEspace.lean\#L226}

\begin{theoreme}
Caractérisation du cercle de diamètre $[AB]$ : le point $M$ y appartient si et seulement si
\[ \overrightarrow{MA} \cdot \overrightarrow{MB} = 0 . \]
\end{theoreme}

\begin{proof}
Posons $u = m - a$ et $v = m - b$, de sorte que
$\overrightarrow{MA}\cdot\overrightarrow{MB} = \langle u, v\rangle$ \textemdash{} les deux
vecteurs $a - m$ et $b - m$ étant les opposés de $u$ et $v$, et le produit scalaire de deux
opposés étant inchangé.

Si $I$ désigne le milieu de $[AB]$, alors
\[ m - I = \tfrac{1}{2}(u + v) , \qquad a - b = v - u , \]
la première égalité parce que $I = \frac{1}{2}(a+b)$. L'appartenance au cercle,
$MI = \frac{AB}{2}$, s'écrit donc
\[ \tfrac{1}{2}\|u + v\| = \tfrac{1}{2}\|v - u\| , \]
c'est-à-dire l'égalité des deux normes. On est ramené au critère d'orthogonalité démontré
plus haut : en développant
\[ \|u+v\|^{2} = \|u\|^{2} + 2\langle u, v\rangle + \|v\|^{2} , \qquad
   \|v-u\|^{2} = \|v\|^{2} - 2\langle v, u\rangle + \|u\|^{2} , \]
et en utilisant la symétrie du produit scalaire, la différence des carrés vaut
$4\langle u, v\rangle$. Les deux normes, positives, sont donc égales si et seulement si
$\langle u, v\rangle = 0$.

\medskip\noindent
C'est le théorème de l'angle inscrit dans un demi-cercle, dû à Thalès, sous sa forme
vectorielle.
\end{proof}
\source{https://github.com/Commutator-IO/learn-lean/blob/main/cours/02-lycee/08-geometrie-vecteurs-produit-scalaire-espace/GeometrieVecteursProduitScalaireEspace.lean\#L241}{GeometrieVecteursProduitScalaireEspace.lean\#L241}

\subsection{Géométrie dans l'espace : droites et plans}

\begin{theoreme}
Si deux points d'une droite appartiennent à un plan, la droite entière y est incluse :
\[ A \in \mathcal{P},\ B \in \mathcal{P} \implies (AB) \subset \mathcal{P} . \]
\end{theoreme}

\begin{proof}
La droite $(AB)$ est le sous-espace affine engendré par $\{A, B\}$, c'est-à-dire le plus
petit qui les contienne. Dire qu'il est inclus dans $\mathcal{P}$ revient donc à dire que
$\{A, B\} \subset \mathcal{P}$, ce qui est l'hypothèse.

\medskip\noindent
\textbf{Ce qui n'est pas démontré.} La classification complète des positions relatives
\textemdash{} deux droites de l'espace sont sécantes, strictement parallèles ou non
coplanaires ; une droite et un plan sont sécants, parallèles ou l'un dans l'autre
\textemdash{} n'est pas formalisée. Ce n'est pas un théorème mais une taxinomie : la
démontrer demanderait de raisonner sur les dimensions des sous-espaces affines et de leurs
intersections, ce que le fichier ne fait pas. Le théorème ci-dessus en est la seule pièce
utile isolée : une droite qui n'est pas dans un plan le coupe en au plus un point, faute de
quoi elle y serait incluse.

\medskip\noindent
Le \textbf{théorème du toit} n'est pas formalisé non plus. Sa démonstration passe par le
même calcul de dimensions : l'intersection de deux plans distincts de l'espace est une
droite parce que $2 + 2 - 3 = 1$, et c'est cette égalité, invisible dans l'énoncé du
lycée, qui porte tout le contenu.
\end{proof}
\source{https://github.com/Commutator-IO/learn-lean/blob/main/cours/02-lycee/08-geometrie-vecteurs-produit-scalaire-espace/GeometrieVecteursProduitScalaireEspace.lean\#L271}{GeometrieVecteursProduitScalaireEspace.lean\#L271}

\begin{theoreme}
Représentation paramétrique d'une droite : la droite passant par $A$ et dirigée par $u$
est l'ensemble des points
\[ M = A + t\,u , \qquad t \in \mathbb{R} . \]
\end{theoreme}

\begin{proof}
Le sous-espace affine passant par $A$ de direction $\mathrm{Vect}(u)$ contient $M$ si et
seulement si le vecteur $M - A$ appartient à cette direction, c'est-à-dire s'il existe $t$
tel que $tu = M - A$. Chaque sens s'obtient en reportant cette égalité dans l'autre : de
$tu = M - A$ on tire $M = A + tu$, et réciproquement.
\end{proof}
\source{https://github.com/Commutator-IO/learn-lean/blob/main/cours/02-lycee/08-geometrie-vecteurs-produit-scalaire-espace/GeometrieVecteursProduitScalaireEspace.lean\#L279}{GeometrieVecteursProduitScalaireEspace.lean\#L279}

\begin{theoreme}
Représentation paramétrique d'un plan : le plan passant par $A$ et dirigé par $u$ et $v$
est l'ensemble des points
\[ M = A + s\,u + t\,v , \qquad (s, t) \in \mathbb{R}^{2} . \]
\end{theoreme}

\begin{proof}
Même démonstration, avec la description de l'espace engendré par deux vecteurs : $M - A$
appartient à $\mathrm{Vect}(u, v)$ si et seulement s'il s'écrit $su + tv$. Le passage entre
$M - A = su + tv$ et $M = A + su + tv$ est un déplacement de terme.
\end{proof}
\source{https://github.com/Commutator-IO/learn-lean/blob/main/cours/02-lycee/08-geometrie-vecteurs-produit-scalaire-espace/GeometrieVecteursProduitScalaireEspace.lean\#L291}{GeometrieVecteursProduitScalaireEspace.lean\#L291}

\begin{theoreme}
Coplanarité : un vecteur $w$ est coplanaire à $u$ et $v$ si et seulement s'il en est
combinaison linéaire,
\[ w \in \mathrm{Vect}(u, v) \iff \exists (s, t) \in \mathbb{R}^{2},\ w = su + tv . \]
\end{theoreme}

\begin{proof}
C'est la description de l'espace engendré par deux vecteurs, citée à la bibliothèque, à
l'ordre près de l'égalité.
\end{proof}
\source{https://github.com/Commutator-IO/learn-lean/blob/main/cours/02-lycee/08-geometrie-vecteurs-produit-scalaire-espace/GeometrieVecteursProduitScalaireEspace.lean\#L306}{GeometrieVecteursProduitScalaireEspace.lean\#L306}

\begin{theoreme}
Base de l'espace : si $u$, $v$, $w$ sont linéairement indépendants dans un espace de
dimension $3$, tout vecteur $x$ s'écrit d'une seule façon
\[ x = \alpha u + \beta v + \gamma w . \]
\end{theoreme}

\begin{proof}
\emph{Existence.} Trois vecteurs indépendants dans un espace de dimension trois en forment
une base \textemdash{} la bibliothèque construit cette base à partir de l'indépendance et
de l'égalité entre le nombre de vecteurs et la dimension. Les coordonnées de $x$ dans cette
base, notées $x_0$, $x_1$, $x_2$, vérifient par construction
$x = x_0 u + x_1 v + x_2 w$, la somme sur les trois indices étant développée.

\emph{Unicité.} Soit $x = \alpha u + \beta v + \gamma w$ une autre décomposition. En
retranchant les deux égalités,
\[ (\alpha - x_0)\,u + (\beta - x_1)\,v + (\gamma - x_2)\,w = 0 . \]
L'indépendance linéaire dit précisément qu'une combinaison nulle a tous ses coefficients
nuls : $\alpha = x_0$, $\beta = x_1$, $\gamma = x_2$.
\end{proof}
\source{https://github.com/Commutator-IO/learn-lean/blob/main/cours/02-lycee/08-geometrie-vecteurs-produit-scalaire-espace/GeometrieVecteursProduitScalaireEspace.lean\#L317}{GeometrieVecteursProduitScalaireEspace.lean\#L317}

\subsection{Produit scalaire dans l'espace : vecteur normal et équation cartésienne}

\begin{theoreme}
Équation cartésienne d'un plan de vecteur normal $\vec n(a, b, c)$ passant par $P$ :
\[ \langle \vec n, \overrightarrow{PM}\rangle = 0
   \iff a x + b y + c z + d = 0 ,
   \qquad d = -\bigl(a x_P + b y_P + c z_P\bigr) . \]
\end{theoreme}

\begin{proof}
Le produit scalaire dans l'espace repéré est la somme des produits des coordonnées :
\[ \langle \vec n, \overrightarrow{PM}\rangle
   = n_0 (m_0 - p_0) + n_1 (m_1 - p_1) + n_2 (m_2 - p_2) . \]
En développant, les termes en $p$ se regroupent en une constante :
\[ = n_0 m_0 + n_1 m_1 + n_2 m_2 - \bigl(n_0 p_0 + n_1 p_1 + n_2 p_2\bigr) , \]
ce qui est exactement $ax + by + cz + d$ avec les notations du programme. L'équivalence
n'est donc qu'une réécriture : c'est le développement qui fait apparaître $d$.
\end{proof}
\source{https://github.com/Commutator-IO/learn-lean/blob/main/cours/02-lycee/08-geometrie-vecteurs-produit-scalaire-espace/GeometrieVecteursProduitScalaireEspace.lean\#L346}{GeometrieVecteursProduitScalaireEspace.lean\#L346}

\begin{theoreme}
Distance d'un point à un plan : si $\vec n \ne \vec 0$ et si $\mathcal{P}$ est le plan
passant par $P$ de vecteur normal $\vec n$, alors tout point $X$ de $\mathcal{P}$ vérifie
\[ MX \ge \frac{\bigl|\langle \vec n, \overrightarrow{PM}\rangle\bigr|}{\|\vec n\|} , \]
et le point
$X_0 = M - \dfrac{\langle \vec n, \overrightarrow{PM}\rangle}{\|\vec n\|^{2}}\,\vec n$
appartient à $\mathcal{P}$ et réalise l'égalité.
\end{theoreme}

\begin{proof}
\emph{La minoration.} Soit $X$ un point du plan, donc
$\langle \vec n, X - P\rangle = 0$. En écrivant $M - X = (M - P) - (X - P)$ et en
développant par linéarité,
\[ \langle \vec n, M - X\rangle = \langle \vec n, M - P\rangle . \]
L'inégalité de Cauchy--Schwarz donne alors
\[ \bigl|\langle \vec n, M - P\rangle\bigr| = \bigl|\langle \vec n, M - X\rangle\bigr|
   \le \|\vec n\|\,\|M - X\| , \]
et l'on divise par $\|\vec n\| > 0$.

\emph{Le cas d'égalité.} Posons
$\lambda = \dfrac{\langle \vec n, M - P\rangle}{\|\vec n\|^{2}}$ et
$X_0 = M - \lambda \vec n$. Alors
\[ \langle \vec n, X_0 - P\rangle
   = \langle \vec n, M - P\rangle - \lambda \langle \vec n, \vec n\rangle
   = \langle \vec n, M - P\rangle - \lambda \|\vec n\|^{2} = 0 , \]
donc $X_0$ est bien dans le plan ; et
\[ MX_0 = \|\lambda \vec n\| = |\lambda|\,\|\vec n\|
   = \frac{\bigl|\langle \vec n, M - P\rangle\bigr|}{\|\vec n\|} . \]
Le point $X_0$ est le projeté orthogonal de $M$ sur $\mathcal{P}$, et la distance du point
au plan est bien le minimum annoncé.

\medskip\noindent
En coordonnées, avec $\vec n(a, b, c)$, on retrouve la formule du programme :
\[ d(M, \mathcal{P})
   = \frac{|a x_M + b y_M + c z_M + d|}{\sqrt{a^{2} + b^{2} + c^{2}}} . \]

\medskip\noindent
\textbf{Ce qui n'est pas démontré.} Les intersections droite-plan et plan-plan ne sont pas
formalisées : comme la classification des positions relatives, elles reposent sur un calcul
de dimensions absent de ce fichier.
\end{proof}
\source{https://github.com/Commutator-IO/learn-lean/blob/main/cours/02-lycee/08-geometrie-vecteurs-produit-scalaire-espace/GeometrieVecteursProduitScalaireEspace.lean\#L357}{GeometrieVecteursProduitScalaireEspace.lean\#L357}

\vfill
\noindent\rule{0.35\linewidth}{0.4pt}\par
\noindent{\footnotesize Leonardo de Moura et Sebastian Ullrich,
\emph{The Lean~4 Theorem Prover and Programming Language},
\emph{Automated Deduction — CADE~28}, Springer, 2021, p.~625--635,
\href{https://doi.org/10.1007/978-3-030-79876-5_37}{doi:10.1007/978-3-030-79876-5\_37}.\par}

\end{document}
